% Originally: content/week-07.tex, \section{Solutions}

\section{Solutions}
\label{sec:ift_solutions}

% Transition: Claude Opus 5 (High)
Not every solution in this chapter is collected here. \Cref{ex:open_mapping_grad_nonzero} is solved
where it appears, on \cpageref{ex:open_mapping_grad_nonzero}.

\begin{exercisesolution}[Solution to \cref{ex:7.1}]
% Generator: Claude Sonnet 5 (Medium)
\begin{enumerate}[label=\textbf{(\alph*)}]
  \item First build a diffeomorphism $\mathbb{R} \to (0,1)$: the function
    \[\psi : \mathbb{R} \to (0,1), \qquad \psi(t) := \frac12+\frac1\pi\arctan(t)\]
    is $C^\infty$, strictly increasing (since $\arctan' > 0$), and bijective onto $(0,1)$ (as
    $\arctan$ is bijective onto $(-\tfrac\pi2,\tfrac\pi2)$), with $C^\infty$ inverse
    $\psi^{-1}(s) = \tan\big(\pi(s-\tfrac12)\big)$. Applying $\psi$ in each coordinate gives the
    diffeomorphism
    \[\phi : \mathbb{R}^2 \to (0,1)\times(0,1), \qquad \phi(x,y) := \Big(\frac12+\frac1\pi\arctan(x),\ \frac12+\frac1\pi\arctan(y)\Big).\]
  \item \textbf{No.} $f$ is a $C^\infty$ bijection of $\mathbb{R}$ with $f^{-1}(x) =
    \sqrt[5]{x}$, but $f'(0) = 5\cdot0^4 = 0$. If $f$ were a diffeomorphism, then by the
    chain rule applied to $f^{-1}\circ f = \id$,
    \[1 = (f^{-1})'(f(x))\cdot f'(x) \quad \text{for all } x,\]
    which is impossible at $x=0$ since $f'(0)=0$. Equivalently, $(f^{-1})'(x) =
    \tfrac15x^{-4/5} \to \infty$ as $x \to 0$, so $f^{-1}$ is not even differentiable at $0$,
    let alone continuously so.
\end{enumerate}
\end{exercisesolution}

\begin{exercisesolution}[Solution of \cref{ex:hessian_gradient_diffeomorphism}]
% Generator: Claude Opus 5 (High)
Write $H := \Hess f(0)$. Everything follows from its spectrum, so compute that first. The matrix is
block diagonal, with the $1\times1$ block $(1)$ and the $2\times2$ block
$\left(\begin{smallmatrix} 0 & 1 \\ 1 & 0\end{smallmatrix}\right)$, whose eigenvalues are $\pm 1$.
Hence $H$ has eigenvalues $1, 1, -1$, and in particular $\det H = -1 \neq 0$: the critical point is
non-degenerate, and $H$ is indefinite.

Since $f$ is smooth with $f(0) = 0$ and $\nabla f(0) = 0$, Taylor's theorem gives
\begin{equation}
\label{eq:degenerate_critical_taylor}
f(x) = \tfrac12 \langle Hx, x\rangle + o(\lvert x\rvert^2)
     = \tfrac12 x_1^2 + x_2x_3 + o(\lvert x\rvert^2)
     \qquad \text{as } x \to 0 .
\end{equation}
\begin{enumerate}[label=\textbf{(\alph*)}]
  \item \textbf{False.} Along the $x_1$-axis, \eqref{eq:degenerate_critical_taylor} gives
    $f(t,0,0) = \tfrac12 t^2 + o(t^2) > 0 = f(0)$ for small $t \neq 0$, so $f$ takes values larger
    than $f(0)$ arbitrarily close to the origin.
  \item \textbf{True.} We have just seen a direction along which $f$ increases. In the plane
    $x_1 = 0$ the quadratic part is $x_2x_3$, so along $x_2 = -x_3 = t$ it reads $-t^2$, and
    $f(0,t,-t) = -t^2 + o(t^2) < 0 = f(0)$. Taking both values above and below $f(0)$ in every
    neighbourhood of the origin is exactly what a saddle point does.
  \item \textbf{False.} The quadratic form is correct only up to a change of coordinates, and the
    statement is made in the original ones. By \eqref{eq:degenerate_critical_taylor} the second-order
    part of $f$ is $\tfrac12 x_1^2 + x_2 x_3$, which is not $x_1^2 - x_2^2 - x_3^2$: evaluating both
    at $(0,1,1)$ gives $1$ and $-2$. What \emph{is} true is that in suitable coordinates the form
    becomes $\tfrac12(u_1^2 + u_2^2 - u_3^2)$, matching the eigenvalues $1,1,-1$ --- note in
    passing that this has signature $(2,1)$, not the $(1,2)$ the stated expression suggests.
  \item \textbf{True.} This is the same Taylor expansion read with a weaker conclusion: the
    right-hand side of \eqref{eq:degenerate_critical_taylor} is bounded by a constant times
    $\lvert x\rvert^2$ near $0$. Vanishing to first order at a point is precisely what $O(\lvert
    x\rvert^2)$ records.
  \item \textbf{True}, and this is where the chapter's theorem enters. The map $\nabla f$ is smooth,
    being the gradient of a smooth function, and its Jacobian at the origin is
    $\Jac(\nabla f)(0) = \Hess f(0) = H$, which we computed to be invertible. By the inverse function
    theorem (\cref{thm:inverse_function_theorem}), $\nabla f$ is therefore a diffeomorphism from some
    ball $B_\varepsilon(0)$ onto an open neighbourhood of $\nabla f(0) = 0$.
\end{enumerate}
\begin{remark}
Item \textbf{(e)} is the useful one to remember: \qt{non-degenerate critical point} means exactly
\qt{$\nabla f$ is a local diffeomorphism there}, so such points are isolated. Item \textbf{(c)} is
the standard trap --- diagonalising the Hessian is a statement about \emph{some} coordinates, and
carrying its conclusion back to the original ones without the change of variables is where the
argument breaks.
\end{remark}
\end{exercisesolution}

\begin{exercisesolution}[Solution of \cref{ex:local_diffeo_parameter_alpha}]
% Generator: Claude Opus 5 (High)
Everything turns on one number:
\[\det\Jac\Phi(0,0) = (\alpha+5)(\alpha-5) - (-3)\cdot 3 = \alpha^2 - 25 + 9 = \alpha^2 - 16 .\]
\begin{enumerate}[label=\textbf{(\alph*)}]
  \item \textbf{False.} The determinant vanishes exactly when $\alpha^2 = 16$, so $\alpha^2 \neq 25$
    is the wrong condition. Take $\alpha = 4$, which satisfies $\alpha^2 = 16 \neq 25$ and so is
    covered by the claim. There $\det\Jac\Phi(0,0) = 0$, and no restriction of $\Phi$ to a
    neighbourhood of the origin can be a diffeomorphism: if $\Phi|_V$ had a $C^1$ inverse $\Psi$,
    the chain rule applied to $\Psi\circ\Phi = \id$ would give
    $\Jac\Psi(\Phi(0,0))\,\Jac\Phi(0,0) = I_2$, forcing $\Jac\Phi(0,0)$ to be invertible.
  \item \textbf{True.} If $\alpha^2 \neq 16$, then $\det\Jac\Phi(0,0) \neq 0$, so
    $D\Phi_{(0,0)}$ is invertible, and $\Phi \in C^1$. The inverse function theorem
    (\cref{thm:inverse_function_theorem}) supplies open sets $V \ni (0,0)$ and $W \ni \Phi(0,0)$
    with $\Phi|_V : V \to W$ a $C^1$-diffeomorphism. Any ball $B_r(0,0) \subseteq V$ then works,
    since a restriction of a diffeomorphism to an open subset is again one onto its image.
  \item \textbf{True.} For $\alpha = 0$ the determinant is $-16 \neq 0$, so \textbf{(b)} applies.
    Fix $V, W$ as there and $r > 0$ with $B_r(0,0) \subseteq V$. A diffeomorphism is in particular
    a homeomorphism onto the open set $W$, so it maps open subsets of $V$ to open subsets of $W$,
    and open subsets of an open set are open in $\mathbb{R}^2$. Hence $\Phi(B_r(0,0))$ is open.
\end{enumerate}
\begin{remark}[Where $25$ comes from, and why it is wrong]
The two thresholds are not a random pair. A reader who computes the determinant as a product of
the diagonal entries and forgets the off-diagonal contribution gets $(\alpha+5)(\alpha-5)$ and
lands on $\alpha^2 = 25$, which is exactly the false condition of \textbf{(a)}. The off-diagonal
pair $-3, 3$ contributes $+9$ and moves the singular set to $\alpha^2 = 16$, that is
$\alpha = \pm4$. Item \textbf{(a)} and item \textbf{(b)} are therefore the same sentence with one
arithmetic slip between them, and reading them as a pair is faster than checking either alone.
\end{remark}
\end{exercisesolution}

\begin{exercisesolution}[Solution of \cref{ex:kobel_ift_single_choice}]
% Generator: Claude Opus 5 (High)
The correct answer is \textbf{(b)}.
\begin{enumerate}[label=\textbf{(\alph*)}]
  \item \textbf{False}, and it fails before any hypothesis on the derivative is examined.
    \Cref{thm:inverse_function_theorem} needs domain and codomain of the \emph{same} dimension,
    because a local inverse has to be defined on a neighbourhood of the image point. For
    $n > 2$ the image of $f$ has empty interior in $\mathbb{R}^n$ --- it is at best a
    $2$-dimensional surface --- so no neighbourhood of $f(0,0)$ in $\mathbb{R}^n$ lies in the
    image, and there is nothing for an inverse to be defined on. Being locally \emph{injective},
    which is what an immersion gives, is a strictly weaker conclusion.
  \item \textbf{True.} For a map $\mathbb{R}^2 \to \mathbb{R}^2$, maximal rank means rank $2$,
    which for a square matrix is invertibility, that is $\det\Jac f(0,0) \neq 0$. That is exactly
    the hypothesis of \cref{thm:inverse_function_theorem}, and $f \in C^2 \subseteq C^1$ supplies
    the regularity.
  \item \textbf{False.} One non-vanishing partial derivative is one column of the Jacobian, and a
    matrix with one non-zero column can still be singular. Take $f(x,y) := (x,0)\transp$: then
    $\partial f/\partial x = (1,0)\transp \neq 0$ everywhere, while
    $\Jac f = \left(\begin{smallmatrix}1&0\\0&0\end{smallmatrix}\right)$ has determinant $0$, and
    $f$ is not injective on any neighbourhood of the origin --- it collapses every vertical
    segment to a point.
  \item \textbf{False}, by the same example with the roles of the variables exchanged: for
    $f(x,y) := (y,0)\transp$ the partial $\partial f/\partial y$ is non-zero and the Jacobian is
    again singular.
\end{enumerate}
\begin{remark}[The same answer if the partials are read component-wise]
Items \textbf{(c)} and \textbf{(d)} are printed with $f$ itself, so
$\partial f/\partial x(0,0) \neq 0$ says a whole column of the Jacobian is non-zero. A reader who
takes $f = (f_1,f_2)\transp$ and reads them as $\partial f_1/\partial x(0,0) \neq 0$ gets a weaker
hypothesis --- a single \emph{entry} is non-zero --- and the same counterexample refutes it. The
question is robust to the ambiguity, which is presumably why it went unnoticed.

What both readings test is the difference between \qt{some derivative does not vanish} and
\qt{the derivative is invertible}. Only the second is a statement about the linear map $Df_{(0,0)}$
as a whole, and only the second is what \cref{rem:inverse_function_theorem_intuition} argues is the
right hypothesis.
\end{remark}
\end{exercisesolution}
